\loai{Tính tần số}
\begin{vidu} %[3P3-12.2B][Loai1] 
	\nguon{CĐ - 2013} Một máy phát điện xoay chiều một pha có phần cảm là rôto gồm 6 cặp cực (6 cực nam và 6 cực bắc). Rôto quay với tốc độ 600 vòng/phút. Suất điện động do máy tạo ra có tần số bằng
	\choice
	{\True 60 Hz}
	{100 Hz}
	{50 Hz}
	{120 Hz}
	\loigiai{
%		\Binhluan{
		\begin{itemize}
			\item Tốc độ 600 vòng/phút $=$ 10 vòng/giây $= n$. 
			\item Ta có: $f= pn=6.10=60\ Hz$. 
		\end{itemize}
%	}
%{
%Chú ý đổi đơn vị $[n]=$vòng/s. p là số "\textbf{cặp} cực" chứ không phải "số cực".
%}
	}
\end{vidu}
\loai{Tính số cặp cực từ}
\begin{vidu} %[3P3-12.2B][Loai2] 
	Một cỗ máy của nhà máy thủy điện Hòa Bình có rôto quay đều với vận tốc 125 vòng/phút. Số cặp cực từ của máy phát điện của tổ máy là
	\choice
	{\True 24}
	{48}
	{125}
	{12}
	\loigiai{
%		\Binhluan{
		\begin{itemize}
			\item Do tốc độ quay của rôto là: $n=125$ (vòng /phút) nên áp dụng công thức: 
			$f=\dfrac{np}{60}\Rightarrow p=\dfrac{60f}{n}$.
			\item Tần số điện xoay chiều của Việt Nam là $f=50\ Hz\Rightarrow $ thay vào ta có $p=24$ cặp cực từ.
		\end{itemize}
%	}
%{
%Bài toán yêu cầu học sinh phải nhớ tần số của mạng điện xoay chiều ở Việt Nam là $f=50\ Hz$.
%}
	}
\end{vidu}
\begin{vidu} %[3P3-12.2B][Loai2] 
	Một máy phát điện có phần ứng gồm 4 cuộn dây giống nhau mắc nối tiếp, mỗi cuộn có 62 vòng. Để tạo ra suất điện động hiệu dụng 220 V thì rôto phải quay với tốc độ 10 vòng/s. Biết từ thông cực đại qua mỗi vòng dây là 4 mWb. Rôto có số cặp cực là
	\choice
	{2}
	{3}
	{4}
	{\True 5}
	\loigiai{
%\Binhluan{
		Ta có: 
		\begin{align*}
		&E_0 = \omega N\Phi_0 \Leftrightarrow \sqrt{2}E=2\pi f N\Phi_0 \Rightarrow f = \dfrac{\sqrt{2}E}{2\pi  N\Phi_0} =\dfrac{220\sqrt{2}}{2\pi. 4.62.4.10^{-3}} = 50\\ 
		\Rightarrow\ &p =\dfrac{f}{n}=\dfrac{50}{10}=5.
		\end{align*}
%}
%{
%4 cuộn dây. Mỗi cuộn có 62 vòng nên tổng số vòng dây là $N=4.62=248$ vòng.
%}
	}
\end{vidu}
\loai{Tính số cực từ}
\begin{vidu} %[3P3-12.2B][Loai3] 
	Máy phát điện xoay chiều một pha sinh ra suất điện động $e = E_0\cos 120\pi t$ V. Rôto là phần cảm và quay với tốc độ 600 vòng/phút. Số cực từ của rôto mắc xen kẽ ở phần cảm là
	\choice
	{\True 12 cực}
	{10 cực}
	{6 cực}
	{24 cực}
	\loigiai{
%\Binhluan{
		\begin{itemize}
		\item Đổi 600 vòng/phút $=$ 10 vòng/s.
		\item Ta có: $\omega  = 120\pi  \Rightarrow f = pn = 60\ Hz \Rightarrow p = 6$ cặp cực từ $\Rightarrow$ 12 cực từ.
	\end{itemize}
%	}
%{
%Đề bài hỏi \textbf{số cực từ} $= 2p$. 
%}
	}
\end{vidu}
\loai{Tính tốc độ quay của rôto}
\begin{vidu} %[3P3-12.2B][Loai4] 
	Một máy phát điện xoay chiều một pha có rôto gồm 4 cặp cực phát ra suất điện động xoay chiều có tần số 60 Hz. Tốc độ quay của rôto là
	\choice
	{450 vòng/phút}
	{1500 vòng/phút}
	{\True 900 vòng/phút}
	{1800 vòng/phút}
	\loigiai{
%		\Binhluan{
		Ta có: $n=\dfrac{f}{p}=15 $ vòng/s $= 900$ vòng/phút.
	%}
	%{Áp dụng theo đúng công thức sẽ ra đơn vị của tốc độ quay rôto là vòng/s}
	}
\end{vidu}
\loai{Đồ thị}
\begin{vidu}%[3P3-12.2K][Loai10]
	immini[thm]{Máy phát điện xoay chiều một pha, nam châm có 10 cặp cực quay với tốc độ n (vòng/phút) tạo ra suất điện động có đồ thị phụ thuộc thời gian như hình vẽ. Tính n.
		\choice
		{50}
		{\True 100}
		{150}
		{200}}{\begin{tikzpicture}[scale=1,thick,>=stealth',x=1.5cm,y=1cm]
		\draw[->] (0,0)--(3.5,0)node[below]{$t\ (ms)$};
		\draw[->] (0,-2)--(0,2)node[right]{$e\ (V)$};
		\foreach  \x  in  {0,0.2,...,3}
		{\draw [dashed,thin] (\x,-1.5) --(\x,1.5);}
		
		\foreach  \A  in  {1.5} % Biên độ
		\foreach  \T  in  {2.4} % Chu kì
		\foreach  \Phi  in  {-30} % Pha ban đầu
		{\draw [red, line width = 1.2pt, domain=0:3.2, samples=500]%
			plot (\x, {(\A)*cos(360/\T*\x+\Phi)});
		}
		\node at (0,0)[left]{$O$};
		\node at (0.1,0)[below right]{$5$};
		\draw[fill=white](0.2,0)circle(1.2pt);
		\end{tikzpicture}}
	\loigiai{
		%\Binhluan{
			\begin{itemize}
				\item Chu kì $T = 12$ ô $= 60$ ms $\Rightarrow f = \dfrac{50}{3}\ Hz$.
				\item Mà $f = \dfrac{np}{60} \Rightarrow n = \dfrac{60f}{p}=\dfrac{\dfrac{60.50}{3}}{10}= 100$ vòng/phút.
			\end{itemize}
		%}
		%{
		%	Chú ý trục của t là ms.
		%}
	}
\end{vidu}
\loai{Tính suất điện động}
\begin{vidu} %[3P3-12.2B][Loai5] 
	Rôto của máy phát điện xoay chiều một pha là một nam châm có 4 cặp cực từ, quay với tốc độ 1500 vòng/phút. Mỗi cuộn dây của phần ứng có 50 vòng dây. Biết ứng với mỗi cực từ lại có một cuộn dây được quấn trên stator, các cuộn dây giống nhau Từ thông cực đại qua mỗi vòng dây là 5 mWb. Suất điện động cảm ứng hiệu dụng do máy tạo ra là
	\choice
	{1256 V}
	{\True 888 V}
	{444 V}
	{628 V}
	\loigiai{
		%\Binhluan{
		\begin{itemize}
			\item Rôto của máy phát có 4 cặp cực từ $\Rightarrow$ phần ứng gồm 8 cuộn dây.
			\item Tốc độ góc $\omega =\dfrac{2\pi np}{60}=200\pi $ rad/s.
			$\Rightarrow E=\dfrac{E_0}{\sqrt{2}}=\dfrac{\omega NBS}{\sqrt{2}}=\dfrac{\omega N\Phi_0}{\sqrt{2}}=888\ V$.
		\end{itemize}
%}
%{
%Theo cấu tạo của máy phát điện xoay chiều một pha thì mỗi cực của nam châm sẽ ứng với một cuộn dây.
%}
	}
\end{vidu}
\loai{Tính từ thông}
\begin{vidu} %[3P3-12.2B][Loai6] 
	Một máy phát điện xoay chiều một pha có rôto 4 cực từ quay đều với tốc độ 1200 vòng/phút, tạo ra suất điện động hiệu dụng 220 V. Biết ứng với mỗi cực từ lại có một cuộn dây được quấn trên stator, các cuộn dây giống nhau, mỗi cuộn có 100 vòng dây. Từ thông cực đại qua mỗi vòng dây là
	\choice
	{\True 3 mWb}
	{$7,7. 10^{-4}$ Wb}
	{1,5 mWb}
	{2,2 mWb}
	\loigiai{
		%\Binhluan{
		\begin{itemize}
			\item Ta có: $n = 20$ vòng/s $\Rightarrow f = pn = 40$ Hz $\Rightarrow \omega  = 2\pi f = 80\pi$ rad/s; [số cuộn dây] $=$ 4.
			\item Suất điện động cực đại: $E_0 = \omega N.\Phi_0 \Rightarrow \Phi_0 \approx 3$ mWb.
		\end{itemize}
	%}
%{
%Đề bài hỏi từ thông cực đại qua \textbf{mỗi vòng dây} là: $\Phi_0=BS$. 
%}
}
\end{vidu}
\loai{Tính số vòng dây ở mỗi cuộn}
\begin{vidu} %[3P3-12.2B][Loai8] 
	\nguon{ĐH - 2011} Một máy phát điện xoay chiều một pha có phần ứng gồm bốn cuộn dây giống nhau mắc nối tiếp. Suất điện động xoay chiều do máy phát sinh ra có tần số 50 Hz và giá trị hiệu dụng $100\sqrt{2}$ V. Từ thông cực đại qua mỗi vòng dây của phần ứng là $\dfrac{5}{\pi }$ mWb. Số vòng dây trong mỗi cuộn dây của phần ứng là
	\choice
	{400 vòng}
	{\True 100 vòng}
	{71 vòng}
	{200 vòng}
	\loigiai{
	%\Binhluan{
		\begin{itemize}
			\item Gọi tổng số vòng dây của máy là N, ta có: 
			\begin{align*}
				E_0 = E\sqrt{2}= 2\pi f.N.\Phi_0 \Rightarrow N = \dfrac{E\sqrt{2}}{2\pi f{\Phi }_{0}}=\dfrac{100.2}{2.\pi .50.\dfrac{{5.10}^{-3}}{\pi }}=400\ \ct{vòng}.
			\end{align*}
			\item Số vòng dây của 1 cuộn (máy có 4 cuộn dây): $N_{1cuon}=\dfrac{N}{4}=100$.
	\end{itemize}
%}
%{
%Đề bài hỏi số vòng dây \textbf{trong mỗi cuộn dây}.
%}
}
\end{vidu}
\loai{Tăng liên tục tốc độ quay rôto}
\begin{vidu}%[3P3-12.2B][Loai2] 
	\nguon{QG - 2020} Một máy phát điện xoay chiều một  pha đang hoạt động. Khi rôto của máy quay đều với tốc độ n vòng/s thì suất điện động do máy tạo ra có tần số 60 Hz. Khi rôto quay đều với tốc độ $1,5n$ vòng/s thì suất điện động do máy tạo ra có tần số là 
	\choice
	{20 Hz}
	{180 Hz}
	{\True 90 Hz}
	{40 Hz}
	\loigiai{
\begin{itemize}
	\item 
	\item 
	\item 
\end{itemize}	
}	
\end{vidu}
\begin{vidu} %[3P3-12.2K][Loai9] 
	Một máy phát điện xoay chiều một pha khi tăng tốc độ quay của rôto thêm 3 vòng/s thì tần số dao động của suất điện động do máy tạo ra tăng từ 50 Hz đến 65 Hz và suất điện động hiệu dụng do máy phát tạo ra tăng thêm 30 V. Nếu tăng tiếp tốc độ thêm 3 vòng/s nữa thì suất điện động của máy phát tạo ra là
	\choice
	{320 V}
	{280 V}
	{240 V}
	{\True 160 V}
	\loigiai{
	%	\Binhluan{
		\begin{itemize}
			\item Ta có: $50\ Hz = pn$ và $65\ Hz = p.\left(n + 3\right) \Rightarrow p = 5$ và $n = 10$ vòng/s.
			\item Mà $E \sim  n \Rightarrow 30\ V \sim$  3 vòng/s $\Rightarrow$ tăng tiếp thêm 3 vòng/s thì tốc độ quay của rôto là $ 10 + 3 + 3 = 16$ vòng/s ứng với suất điện động là 160 V.
		\end{itemize}
%	}
%{
%Có thể dùng bảng để kiểm soát số liệu tốt hơn.
%}
%	\begin{bang}[tabularx={X|X|X|X}]{}
%		 & n & $f=np$ &  $E=2\pi.f.NBS$ \\\hline
%		\textbf{Lần 1} & n & $50=np$ & $E_1=100\pi.NBS$ \\\hline
%		\textbf{Lần 2} & $n+3$ & $65=\left(n+3\right).p$ & $E_2=E_1+30=130\pi .NBS$ \\\hline
%		\textbf{Lần 3} & $n+6$ & $f_3=\left(n+6\right).p$ & $E_3=2\pi f_3.NBS$ \\\hline
%	\end{bang}
	}
\end{vidu}

%\begin{bang}[tabularx={X|X|X|X|X|X|X|X}]{BẢNG ĐÁP ÁN}
%	1 & 2 & 3 & 4 & 5 & 6 \\\hline
%	7 & 8 & 9 & 10 & 11 & 12 \\\hline
%	13 & 14 & 15 & 16 & 17 & 18 
%\end{bang}
\loai{Hai máy phát}
\begin{vidu} %[3P3-12.2G][Loai7] 
	Hai máy phát điện xoay chiều một pha đang hoạt động bình thường và tạo ra hai suất điện động có cùng tần số f. Rôto của máy thứ nhất có $p_1$ cặp cực và quay với tốc độ $n_1 = 1800$ vòng/phút. Rôto của máy thứ hai có $p_2 = 4$ cặp cực và quay với tốc độ $n_2$. Biết $n_2$ có giá trị trong khoảng từ 12 vòng/giây đến 18 vòng/giây. Giá trị của f là
	\choice
	{54 Hz}
	{\True 60 Hz}
	{48 Hz}
	{50 Hz}
	\loigiai{
%		\Binhluan{
			\begin{itemize}
				\item Ta có: $n_1 = 1800$ vòng/phút $=$ 30 vòng/s.
				\item Lại có: $f = n_1p_1 = n_2p_2 \Rightarrow n_2=\dfrac{p_1}{p_2}n_1=\dfrac{p_1}{4}.30=7,5p_1 $.
				\item Thế vào điều kiện 12 vòng/giây $< n_2 < 18$ vòng/giây.
				\item $12 < 7,5p_1 < 18 \Rightarrow 1,6 < p_1 < 2,4 \Rightarrow p_1=2$ (cặp cực) $\Rightarrow f=n_1p_1=30.2=60\ Hz $. 
			\end{itemize}
%		}
%		{
%			Do p là số cặp cực nên $p \in \mathbb{N^*}$. Ta tìm cách khai triển $n_2$ theo p rồi kẹp điều kiện và tìm nghiệm nguyên tương ứng. 
%		}
	}
\end{vidu}
